\PassOptionsToPackage{unicode}{hyperref} 
\documentclass{beamer}

% --- Налаштування мови та шрифтів ---
\usepackage[T2A]{fontenc}
\usepackage[utf8]{inputenc}
\usepackage[ukrainian]{babel}
\usepackage{graphicx}

% Фікс для математичних шрифтів та чіткості
\usefonttheme[onlymath]{serif} 

% --- ТЕМА ТА КОЛЬОРИ ---
\usetheme{CambridgeUS}
\usecolortheme{dolphin} 

% --- НАЛАШТУВАННЯ ЯСКРАВОСТІ РАМОК (БЛОКІВ) ---
\setbeamercolor{block title}{fg=white, bg=blue!80!black}
\setbeamercolor{block body}{fg=black, bg=blue!10}
\setbeamercolor{block title example}{fg=white, bg=green!50!black}
\setbeamercolor{block body example}{fg=black, bg=green!10}
\setbeamercolor{block title alert}{fg=white, bg=red!70!black}
\setbeamercolor{block body alert}{fg=black, bg=red!10}

\setbeamercolor{normal text}{fg=black}
\setbeamerfont{block title}{series=\bfseries}

% --- НАЛАШТУВАННЯ ПІДПИСІВ ---
\setbeamertemplate{caption}[numbered] 
\setbeamertemplate{caption label separator}{. } 
\addto\captionsukrainian{\renewcommand{\figurename}{Рис.}}

% --- ДАНІ ТИТУЛЬНОГО СЛАЙДА 
% У квадратних дужках [...] пишіть коротко, у фігурних {...} — повно
\title[Теорема Піфагора]{Застосування теореми Піфагора}
\subtitle{Геометрія та її практичне значення}
\author[Прізвище І.П.]{ Прізвище Ім'я По батькові}
\institute[ФМФ]{Факультет математики, фізики та комп'ютерних наук}
\date{\today}

\begin{document}

% 1. ТИТУЛЬНИЙ АРКУШ
\begin{frame}
    \titlepage
\end{frame}

% 2. ЗМІСТ
\begin{frame}{План презентації}
    \tableofcontents
\end{frame}

% 3. СПИСКИ (Поступова поява)
\section{Історична довідка}
\begin{frame}{Хто такий Піфагор?}
    \begin{itemize}
        \item<1-> Давньогрецький філософ та математик.
        \item<2-> Засновник релігійно-філософської школи піфагорійців.
        \item<3-> \textbf{Основний здобуток:} Доведення знаменитої теореми.
    \end{itemize}
\end{frame}

% 4. МАТЕМАТИКА ТА ЯСКРАВІ РАМКИ
\section{Теоретичні основи}
\begin{frame}{Формулювання теореми}
    \begin{block}{Визначення}
        У прямокутному трикутнику квадрат гіпотенузи дорівнює сумі квадратів катетів.
    \end{block}

    \begin{exampleblock}{Основна формула}
        \[ a^2 + b^2 = c^2 \]
    \end{exampleblock}
    
    \begin{alertblock}{Увага!}
        Тільки для трикутників з прямим кутом ($90^\circ$).
    \end{alertblock}
\end{frame}

% 5. ГРАФІКА ТА КОЛОНКИ
\section{Візуалізація}
\begin{frame}{Графічна інтерпретація}
    \begin{columns}
        \begin{column}{0.5\textwidth}
            \textbf{Єгипетський трикутник:}
            \begin{itemize}
                \item $a = 3, b = 4, c = 5$
                \item<2-> $9 + 16 = 25$
            \end{itemize}
        \end{column}
        
        \begin{column}{0.5\textwidth}
            \centering
            \begin{figure}
                \includegraphics[width=0.8\textwidth]{example-image}
                \caption{Схема трикутника}
            \end{figure}
        \end{column}
    \end{columns}
\end{frame}

% 6. ВИСНОВОК
\section{Висновки}
\begin{frame}{Практичне значення}
    \begin{enumerate}
        \item Будівництво та архітектура.
        \item Навігація та картографія.
        \item Астрономія.
    \end{enumerate}
    \vspace{0.5cm}
    \centering
    \textbf{\Large Дякую за увагу!}
\end{frame}

\end{document}
